\documentclass[a5paper,10pt]{article}

% https://github.com/InDevRus/filippov-solutions
% Нумерация страниц
\pagenumbering{gobble}

% Настройка полей
\usepackage{geometry}
\geometry{left=1cm}
\geometry{right=1cm}
\geometry{top=1cm}
\geometry{bottom=1.5cm}

% Вставка таблицы
\usepackage{array}
\usepackage{tabularx}

% Инструменты выравнивания формул
\usepackage{amsmath}
\usepackage{mathtools}

% Диакритика в математике
\usepackage{amsfonts}

% Вычисления размеров на ходу
\usepackage{calc}

% Удобные записи производных
\usepackage[italic=true]{derivative}

% Настройки сносок
\usepackage{enotez}

% Длинные знаки векторов
\usepackage{anyfontsize}
\usepackage[b]{esvect}

% Вставка картинок
\usepackage{graphicx}

% Вставка красной строки
\usepackage{indentfirst}
\setlength{\parindent}{1em}

% Условия в командах
\usepackage{xifthen}

% Луа-скрипты
\usepackage{luacode}

% Настройка команд отдельно в математическом и текстовом режимах
\usepackage{mathcommand}

% Для повторения LaTeX кода
\usepackage{multido}

% Конфигурация отступов формул
\delimitershortfall-1sp
\usepackage{mleftright}
\mleftright

% Растягивание объектов
\usepackage{relsize}

% Обтекание текстом
\usepackage{wrapfig}
\usepackage{subcaption}
\usepackage{floatflt}

% Поддержка ввода кириллицы
\usepackage[english, russian]{babel}

% Настройка корневой позиции для компилятора
\usepackage{import}
\providecommand*{\ProjectRootPath}{..}

\import{\ProjectRootPath/preambles/macros/}{theorems}

\import{\ProjectRootPath/preambles/fonts}{font_setup}

% Знак градуса
\usepackage{siunitx}

\import{\ProjectRootPath/preambles/macros/}{operators}
\import{\ProjectRootPath/preambles/macros/}{redefinitions}

\import{\ProjectRootPath/preambles/fonts}{letters_setup}
\import{\ProjectRootPath/preambles/fonts/adjustments}{custom_superscripts}

\import{\ProjectRootPath/preambles/custom}{formatting}
\import{\ProjectRootPath/preambles/custom}{frames}
\import{\ProjectRootPath/preambles/custom}{list_setup}

% TODO: перевести команды на современный стиль объявления

\begin{document}

$\tsys{& \dot{x} = -x + 2y - \pow{y}{3} \\& \dot{y} = x - 2y}$.
Возьмём функцию Ляпунова 
$v\br{x; y} = \dfrac {1} {2} \pow{x}{2} + \pow{y}{2}$.
$\der{v}{t} \bigg| 
    = -\pow{x}{2} + 2xy - \linebreak - x\pow{y}{3} + 2xy - 4\pow{y}{2} 
    = -\pow{x}{2} + 4xy - 4\pow{y}{2} - x\pow{y}{3} = -\br{x + 2y}^2 - x\pow{y}{3}$.

Рассматривать точки $\br{x; y}$ будем только в окрестности 
$\vbr{X} = \sqrt{\pow{x}{2} + \pow{y}{2}} < \sub{\epsilon}{0} = 2$.

В области $xy \ge 0$ имеем 
$x\pow{y}{3} = xy \cdot \pow{y}{2} \ge 0$,
и тогда $\der{v}{t} \bigg| \le 0$, причём $\der{v}{t} \bigg| = 0$
тогда и только тогда, когда одновременно $x + 2y = 0$ и $x\pow{y}{3} = 0$.
Последнее возможно только в начале координат. 
Итак, в этой области можно утверждать $\der{v}{t} \bigg| < 0$.

С другой стороны, в области $xy < 0$ получаем 
$\der{v}{t} \bigg| = -\pow{x}{2} - 4\pow{y}{2} + xy \br{4 - \pow{y}{2}}$.
В описанной выше окрестности 
$\vbr{\y} \le \sqrt{\pow{x}{2} + \pow{y}{2}} < \sub{\epsilon}{0} = 2$,
то есть $\pow{y}{2} < 4$ и $4 - \pow{y}{2} > 0$.
Следовательно, 
$xy \br{4 - \pow{y}{2}} < 0$ и $\der{v}{t} \bigg| < 0$,
причём равенство нулю невозможно.

Итак, при всех возможных $\br{x; y}$ внутри окрестности радиуса 
$\sub{\epsilon}{0}$ получаем $\der{v}{t} \bigg| < 0$.
Поскольку $\der{v}{t} \bigg|$ непрерывна, можно взять 
$w\br{x; y} = -\der{v}{t} \bigg|$. 
Теорема Ляпунова гарантирует асимптотическую устойчивость нулевого решения.

\end{document}
